\documentclass{article}
\usepackage{geometry}
\geometry{a4paper,scale=0.8}
\usepackage[utf8]{inputenc}

\usepackage{amsmath, amssymb}
\usepackage{amsthm}
\usepackage{enumitem}

\newtheoremstyle{breakstyle}
    {5pt}
    {10pt}
    {\normalfont}
    {0pt}
    {\bfseries}
    {\newline\indent}
    {0pt}
    {\thmname{#1}\thmnumber{ #2}\thmnote{ \textbf{[#3]}}}
\theoremstyle{breakstyle}
\newtheorem{exercise}{Exercise}[section]
\newtheorem{remark}{Remark}[section]

\title{Quadratic and Cyclotomic Fields}
\author{Flandre}
\date{\today}

\begin{document}
\maketitle

\setcounter{section}{3}

\begin{exercise}
    \leavevmode \vspace{-\baselineskip}
    \begin{enumerate}[label=\alph*)]
        \item From \emph{Theorem 3.3} we know that $\mathbb{Q} \left( \sqrt{3} \right)$ has integral bases $\left\{ 1, \sqrt{3} \right\}$ and discriminant $12$.
        
        \item From \emph{Theorem 3.3} we know that $\mathbb{Q} \left( \sqrt{- 7} \right)$ has integral bases $\left\{ 1, \frac{1}{2} + \frac{i}{2} \sqrt{7} \right\}$ and discriminant $- 7$.
        
        \item From \emph{Theorem 3.3} we know that $\mathbb{Q} \left( \sqrt{11} \right)$ has integral bases $\left\{ 1, \sqrt{11} \right\}$ and discriminant $44$.
        
        \item From \emph{Theorem 3.3} we know that $\mathbb{Q} \left( \sqrt{- 11} \right)$ has integral bases $\left\{ 1, \frac{1}{2} + \frac{i}{2} \sqrt{11} \right\}$ and discriminant $- 11$.
        
        \item From \emph{Theorem 3.3} we know that $\mathbb{Q} \left( \sqrt{6} \right)$ has integral bases $\left\{ 1, \sqrt{6} \right\}$ and discriminant $24$.
        
        \item From \emph{Theorem 3.3} we know that $\mathbb{Q} \left( \sqrt{- 6} \right)$ has integral bases $\left\{ 1, \sqrt{6} i \right\}$ and discriminant $- 6$.
    \end{enumerate}
\end{exercise}

\begin{exercise}
    \leavevmode \vspace{-\baselineskip}
    \begin{enumerate}[label=\alph*)]
        \item $N_K (\zeta^2) = 1$
        
        $T_K (\zeta^2) = - 1$
        
        \item $N_K (\zeta + \zeta^2) = \prod_{k = 1}^4 (\zeta^k + \zeta^{2 k}) = 1$
        
        $T_K (\zeta + \zeta^2) = T_K (\zeta) + T_K (\zeta^2) = - 2$
        
        \item $N_K (1 + \zeta + \zeta^2 + \zeta^3 + \zeta^4) = N_K (0) = 0$
        
        $T_K (1 + \zeta + \zeta^2 + \zeta^3 + \zeta^4) = T_K (0) = 0$
    \end{enumerate}
\end{exercise}

\begin{exercise}
    \leavevmode \vspace{-\baselineskip}
    \begin{enumerate}[label=\roman*.]
        \item When $N_K (\alpha) = \pm 1$, we have $\frac{\pm 1}{\alpha} = \frac{N_K (\alpha)}{\alpha} \in \mathbb{Z} (\zeta)$, so $\alpha$ must have a integer inverse.
        
        \item When $\alpha$ is a unit in $\mathbb{Z} (\zeta)$,
    \end{enumerate}
\end{exercise}

\begin{exercise}
\end{exercise}

\begin{exercise}
\end{exercise}

\begin{exercise}
\end{exercise}

\begin{exercise}
\end{exercise}

\begin{exercise}
\end{exercise}

\begin{exercise}
\end{exercise}

\begin{exercise}
\end{exercise}

\begin{exercise}
\end{exercise}

\end{document}
